\documentclass[]{nenkai2019}  
% Use packages---------------------------------------------
\usepackage[dvips]{graphicx}
\usepackage{amsmath}
%\usepackage{url}
%\usepackage{here}
%\usepackage{hyperref}
%\usepackage[margin=5pt]{subcaption}
\usepackage{color}
\usepackage{ulem} % for the line over texts 
\usepackage{epstopdf}%This line makes .eps figures into .pdf - please comment out if not required.
\definecolor{darkgreen}{RGB}{0,129,0}
\usepackage{bm}
% ------------------------------------------------ %
\def\notyet#1{{\color{blue}#1}} % not yet done 
\def\rev#1{{\color{red}#1}} % means should be reviewed again  
\def\you#1{{\color{darkgreen}#1}} % comments by you

% ------------------------------------------------ %
\title{日本語タイトル\textless14 ptゴシック，Arial\textgreater}

\etitle{English Title \textless12 pt Arial\textgreater}

\author{○著者1, 所属機関略称, 所属機関住所, E-mail: \textless 10 pt 明朝体, Times\textgreater\\
        　著者2, 所属機関略称, 所属機関住所, E-mail:}

\eauthor{Author 1, Affiliation, Address \textless 10 pt Times\textgreater\\
Author 2, Affiliation, Address}

\abstract{This is a simple example of how to prepare the paper for JSFM Meeting. 
The headings should appear as above. The instruction is written in the main body. 
This is a simple example of how to prepare the paper for JSFM Meeting. 
The headings should appear as above. 
The instruction is written in the main body. 
This is a simple example of how to prepare the paper for JSFM Meeting. 
The headings should appear as above. \textless Abstract 100-150 words, 9 pt Times\textgreater
}

\papernum{}


\begin{document}
\maketitle
\section{提出物\textless 9 ptゴシック，Arial\textgreater}
提出物は以下の1点のみです．
\textless 9 pt 明朝体，Times\textgreater
%
\begin{enumerate}
 \item[(1)] 講演原稿（\underline{PDF}，CD-ROMへの掲載）： A4版2段組，
 10枚以下（標準\underline{2$\sim$4枚}），\underline{カラー可}
\end{enumerate}
%
ただしPDFのファイルサイズは\underline{3.5 MB以内}とします．
この容量を超える場合には，実行委員会までご相談下さい．

PDFファイルは，以下の手順で作成して下さい．
まず，Word, TeX等の適当なソフトウェアを用いて図表すべてを貼り込んだ原稿を作成し，その後，Adobe社製Acrobat等を用いてAcrobat 5.0以上互換フォーマットでPDFファイルに変換して下さい（PDFファイルの作成が困難な方は，実行委員会までE-mailでご連絡下さい）．
作成したPDFファイルを年会ホームページ（http://www2.nagare.or.jp/nenkai2017/）の記載にしたがってアップロードすれば，講演原稿の提出は完了です．
なお，実行委員会では，例年ご投稿いただいたPDFファイルの確認作業に多くの時間を要しております．
原稿を投稿される際は文字化けなどが発生していないかよくご確認いただき，原稿の修正を折り返し依頼する必要がないようご協力をお願い致します．
お送りいただいた講演原稿はUSBメモリに収録し，当日受付にて配布します．なお，講演原稿のアップロード時にUSB収録へのご同意をお願いしております．著作権の譲渡はございません．


\section{原稿用紙}
A4版用紙の上下に20 mm，左右に15 mmの余白をとり，本文は原則として9ポイントの文字を使用して印字して下さい．
また，表題，著者名，所属，英文要旨の部分を除いて2段組みで作成して下さい．
英文原稿も受け付けますが，その場合でも2段組みと上下左右の余白を守って下さい．
ページ番号は，フッターに1から順に付けて下さい．

\section{体裁}
論文は体裁は以下の通りです．
\setlength{\leftmargini}{1zw}
\begin{list}{・}{}
 \item 和文題目：14ポイント・ゴシック体, Arialフォントで用紙中央に印字する．
       なお，表題，著者名の変更はできません．
 \item 英文題目：12ポイント・Arial, Helvetica, cm（bold）フォントで用紙中央に印字する．
       Main Wordsの最初の文字のみ大文字とする．
 \item 和文著者名：10ポイント・明朝体で英文題目との間を1行空け，著者氏名と所属機関名略称，住所，（可能ならば）E-mailアドレスを書く．また，講演者の名前の前に○を付ける．
 \item 英文著者名：10ポイント・Times, Times New Roman, cmフォントを用い，英文で著者名，所属機関名，住所を書く．
 \item 英文要旨：9ポイント・Times, Times New Roman, cmフォントを用い，英文著者名との間を1行空け，100-150 words程度で幅150 mmに収まるように記載する．
 \item 本文：英文要旨との間を1行空け，9ポイントの明朝体を使用して記載する．
 \item 章及び節のタイトルにはゴシック体／Arialフォントを用いる．
 \item 図表：鮮明かつ適当な大きさのものを，Fig. 1のように本文中の参照する文章近くの適当な位置に貼付する．図表中の文字及び表題は英文とする．
 \item 文献：例えば，この~\cite{arakawa}ように引用し，末尾にまとめる~\cite{arakawa2}．
\end{list}
\setlength{\leftmargini}{2zw}

\vspace{1em}
\begin{figure}[h]
\begin{center}
\includegraphics[scale=0.25]{fig1.eps}
\end{center}
 \caption{A sample of figure}
 \label{sample}
\end{figure}


\section{原稿提出期限}

原稿の提出締切は{\bfseries 20XX年7月XX日}です．締切を厳守して下さい．期限を過ぎますとキャンセルとみなされる場合があります．


\section{原稿提出先}
\begin{enumerate}
 \item[(1)] 講演原稿のPDFファイルは，年会ホームページ（http://www2.nagare.or.jp/nenkai20xx/）の記載にしたがって提出して下さい． 
 \item[(2)] ご不明な点などは，実行委員会宛にE-mailにてお問い合わせ下さい．\\　日本流体力学会年会20xx実行委員会\\　 
 \item[(3)] このoverleafテンプレートは流体力学会LaTeXテンプレートを用いて作成しました. E-mail: asekimoto@cheng.es.osaka-u.ac.jp 
\end{enumerate}

\begin{thebibliography}{99}
\bibitem{arakawa} 荒川，谷口, "論文の書式について," 日本流体力学会年会2003, 1 (2003), pp. 1-1.
\bibitem{arakawa2} Arakawa, C. and Taniguchi, N., "How to prepare the paper," Proc. 17th CFD Symp., 1 (2003), pp. 1-1.
\end{thebibliography}

\end{document}